\subsection{RQ6: How strong are \toolname's specifications versus the baselines?}
\label{subsec:rq6-strength}

\rev{Validity certifies that a specification is discharged by its verifier, but it does not rank how informative it is. We therefore compare accepted input domains (preconditions), normal-return guarantees (postconditions), and predicates at corresponding loop points (loop invariants) independently.}

\smallskip
\noindent\rev{\textbf{LLM judge and expert review.}}
\rev{Separate judge calls compare preconditions by accepted-domain inclusion, postconditions by bidirectional implication over their common input domain, and invariants at corresponding loop points. In Fig.~\ref{fig:fc_vs_judge_strength}, a direction means broader input acceptance or a stronger predicate, respectively; there is no aggregate ranking. Missing clauses mean \texttt{true}; unestablished implications count as Incomp. Expert review corrected 32 verdicts. \appref{app:strength} reports these corrections, prompts, reviewer details, and an A/B-order audit.}

\rev{Fig.~\ref{fig:fc_vs_judge_strength} shows stronger \toolname\ postconditions and loop invariants. Whereas \autospec\ accumulates sampled, verifier-accepted clauses, including redundant or tautological ones, \toolname\ supplies local symbolic states to guide path-sensitive and inductive predicates.}

\rev{\autospec\ more often declares broader input domains, commonly through a missing or \texttt{true} precondition; \toolname\ explicitly records runtime and memory safety requirements. This measures declared permissiveness, whose value depends on whether the additional inputs have meaningful valid executions.}

\rev{The \specgen\ comparison appears in \appref{app:strength} as cross-language context; the verifier stacks differ.}

\smallskip
\noindent\rev{\textbf{Frama-C/WP cross-check.}}
\rev{Separate Frama-C/WP harnesses cross-check 698 pre/postcondition pairs
and 540 invariant pairs using Frama-C 27.1, Alt-Ergo and Z3, with five
seconds per goal. In \toolname/baseline/equal/incomparable order, the counts
for preconditions, postconditions, and invariants are 90/281/302/25,
297/7/310/84, and 56/47/22/415, respectively.}

\rev{WP assigns a direction only when it proves every corresponding
obligation. An unproved implication enters Incomp whether it is false or true
but beyond the prover; unsupported harnesses and unresolved loop
correspondence also enter Incomp. Timeouts and automation limits therefore
leave many invariants unordered. The judge followed by expert review can
inspect semantic relations in these unresolved pairs, while WP provides an
independent formal check. Fig.~\ref{fig:fc_vs_judge_strength} confirms the
precondition and postcondition trends, but WP orders fewer invariants.}

\begin{figure}[t]
\centering
\rev{\footnotesize \tikz\filldraw[fill=figcontractfill,draw=figcontractstroke] (0,0) rectangle (.16,.16); \toolname\ side\quad
\tikz\filldraw[fill=figgrayfill,draw=figgrayframe] (0,0) rectangle (.16,.16); \autospec\ side\\[-1pt]
\tikz\filldraw[fill=figloopfill,draw=figloopstroke] (0,0) rectangle (.16,.16); Equal\quad
\tikz\filldraw[fill=figgray!10,draw=figgrayframe] (0,0) rectangle (.16,.16); Incomparable}

\newcommand{\strengthbar}[6]{%
  \node[anchor=east,text=\revcolor,font=\footnotesize] at (-2,#1) {#2};
  \path[draw=figcontractstroke,line width=.4pt,fill=figcontractfill] (0,#1-.24) rectangle (#3,#1+.24);
  \path[draw=figgrayframe,line width=.4pt,fill=figgrayfill] (#3,#1-.24) rectangle (#4,#1+.24);
  \path[draw=figloopstroke,line width=.4pt,fill=figloopfill] (#4,#1-.24) rectangle (#5,#1+.24);
  \path[draw=figgrayframe,line width=.4pt,fill=figgray!10] (#5,#1-.24) rectangle (100,#1+.24);
  \node[anchor=west,text=\revcolor,font=\footnotesize] at (102,#1) {#6};}
\begin{tikzpicture}[x=.45mm,y=.46cm,every node/.style={font=\footnotesize}]
  \path[use as bounding box] (-45,.05) rectangle (147,8.45);
  \foreach \x in {0,25,50,75,100} {
    \draw[figgrayframe!70,densely dotted] (\x,.65) -- (\x,7.75);
    \node[anchor=north,text=\revcolor,font=\footnotesize] at (\x,.55) {\x\%};}

  \node[anchor=west,text=\revcolor,font=\footnotesize\bfseries] at (0,8.08) {(a) Preconditions ($N=698$)};
  \strengthbar{7.25}{Judge + expert}{13.75}{55.30}{99.86}{96/290/311/1}
  \strengthbar{6.30}{WP}{12.89}{53.15}{96.42}{90/281/302/25}

  \node[anchor=west,text=\revcolor,font=\footnotesize\bfseries] at (0,5.38) {(b) Postconditions ($N=698$)};
  \strengthbar{4.55}{Judge + expert}{49.00}{51.58}{97.42}{342/18/320/18}
  \strengthbar{3.60}{WP}{42.55}{43.55}{87.97}{297/7/310/84}

  \node[anchor=west,text=\revcolor,font=\footnotesize\bfseries] at (0,2.68) {(c) Loop invariants ($N=540$)};
  \strengthbar{1.85}{Judge + expert}{22.22}{31.85}{32.78}{120/52/5/363}
  \strengthbar{.90}{WP}{10.37}{19.07}{23.15}{56/47/22/415}
\end{tikzpicture}
\caption{\rev{Dimension-wise relations for \autospec\ versus \toolname\ over identical data sets. A directional segment means a more permissive accepted domain for preconditions and a logically stronger predicate for postconditions or invariants. ``Judge + expert'' is the corrected judge result; WP checks the same implications. Right-hand counts follow the legend order.}}
\label{fig:fc_vs_judge_strength}
\end{figure}
\let\strengthbar\relax

\begin{finding}
\rev{\textbf{Finding (RQ6).} Symbolic context changes where specification strength appears: \toolname\ yields more informative postconditions and stronger loop invariants than \autospec\ in the shared C/ACSL setting, while \autospec's broader declared input domains mainly reflect safety assumptions that are not written explicitly. The cross-language comparison against \specgen\ (Java/JML) is reported in \appref{app:strength}; no single total ordering applies across all dimensions and languages.}
\end{finding}
